\section{Rocq: reuse an existing loop theorem, and separate integration from novelty}
At the review pin \texttt{cbf94cb552a3b49509e0eee80e7455d4eb5974c6}, the maintained
\texttt{ProgLibPerm.v} and \texttt{ProgLib219.v} still share an identical 753-line,
36,213-byte proof body after the identifier replacement
\texttt{dag219}\,$\mapsto$\,\texttt{dagPerm}. However, the previously proposed
89-node closed-slice transport is \emph{already implemented in archived prototypes}:
\texttt{common-macros/rocq} and \texttt{collection-first/code/rocq/macro-slice}
report Coq~8.18 component checks. Their production integration remains open.
In particular, the shared state-record identity, actual-program structural utilities,
parameterized contract wrappers, producer route and Rocq~9.1 endpoint rebuild must
still be reconciled. The transport idea itself is not a new recommendation here.

A smaller, genuinely additional simplification is available immediately.
\texttt{ProgLibPerm.z465} (lines 670--678) reproves a register-clearing loop by
induction even though \texttt{sem\_while} and \texttt{while\_tac} already prove that
same pattern in \texttt{w199}. The actual DAG rows are
\[
  199\mapsto(\mathrm{EDecnz}\ 5,\mathrm{EContinue}\ 1),\qquad
  465\mapsto(\mathrm{EDecnz}\ 5,\mathrm{EContinue}\ 1).
\]
Keep the existing public statement
\begin{verbatim}
Lemma z465 a : forall m, get m 5 = a ->
  sem (ENode 465) m XNormal (set m 5 0).
\end{verbatim}
and replace only its proof by
\begin{verbatim}
Proof.
  intros m _. revert m.
  while_tac (fun m : St => m) (fun m : St => set m 5 0).
Qed.
\end{verbatim}
This applies the generic loop theorem directly at node 465; it does not apply
the node-199 contract at the wrong node. The old value parameter and premise are
retained for compatibility, although the generic theorem does not need them.

\paragraph{All obligations.}
Instantiate \texttt{sem\_while} with $n=465$, $r=5$, body
$\mathrm{EContinue}\ 1$, $g(m)=m$, and $P(m)=\mathrm{set}\ m\ 5\ 0$.
The six obligations are the literal DAG equation, $5\leq8$, the existing
\texttt{sem\_cont} contract at continue level 1, preservation of register~5 by
the identity function, $P(m)=m$ when register~5 is zero, and
$P(m)=P(g(\mathrm{set}\ m\ 5\ v))$ when that register is $v+1$.
The final two are finite-record field calculations. No functional extensionality,
new induction, new axiom, or change to the machine is required. The continue level
is essential; replacing it by 2 is rejected by the static negative control.

\paragraph{Measured scope and validation.}
The declaration changes from 9 lines / 475 bytes to 5 lines / 171 bytes:
a reduction of 4 lines and 304 bytes per declaration, and one fewer bespoke
induction. These are static text counts, not measured endpoint proof-term or
compile-time savings. The exact old declaration also occurs in the
\texttt{NEW\_EQ} template of \texttt{regenerate218.py}; that is the producer to
change before regenerating the 218, 217 and 215 libraries. The included patch
is producer-only and deliberately leaves archived evidence untouched.
The proposed proof and typed comparison client are \textbf{not compiled} in this
review. Adoption requires the production toolchain, unaffected \texttt{r466}
clients, assumption reports, regenerated-source consistency, unchanged literal
tables, and the sparse/dense K215 endpoint checks.

The independently written \texttt{audit\_reuse.py} verifies the 15 saved inputs
against their Git blob hashes and produces \texttt{REUSE\_AUDIT.json}. It treats
upstream source only as data and executes no downloaded source or producer.
